\documentclass[unnumsec,webpdf,contemporary,large]{oup-authoring-template}%

\usepackage{booktabs}
\usepackage{microtype}
\usepackage{amsmath}

% Typesetting fixes for the narrow two-column measure. microtype alone does not clear these:
% the overruns come from unbreakable monospace tokens in the serialization section
% (\texttt{json/<entry_id>.json}, \texttt{pdb/<entry_id>.pdb}), which carry no legal breakpoint,
% so TeX sets them past the right margin rather than hyphenating. Measured on this manuscript:
% 14 overfull hboxes before, the worst 60.8pt (~21mm) into the margin.
%
% [htt] lets \texttt material hyphenate, which is what gives those tokens a breakpoint at all.
\usepackage[htt]{hyphenat}
% A last-resort stretch TeX may use only when a line would otherwise overrun. This is the
% conventional remedy for a narrow measure and changes no content; without it the remaining
% long-URL and long-citation lines still overrun even once \texttt can break.
\emergencystretch=3em

\newcommand{\societylogo}{}
\graphicspath{{figures/}}

\begin{document}

\journaltitle{Nucleic Acids Research}
\DOI{}
\copyrightyear{2026}
\pubyear{2026}
\vol{}
\issue{Database issue}
\access{}
\appnotes{Database Issue}

\firstpage{1}

\title[sAMPle: A Continuously Updatable Bioactive Peptide Knowledgebase]{sAMPle: A Continuously Updatable Knowledgebase, Multi-Activity Predictor, and Homology-Aware Benchmarking Platform for Bioactive Peptides}

\author[1]{Daniela Peres Martinez\ORCID{0000-0001-5457-5385}}
\author[1]{Isabel Ladeira Pereira\ORCID{0000-0002-5440-0953}}
\author[1]{Rafaella Sinnott Dias\ORCID{0009-0003-8010-1919}}
\author[1]{Luiza Perez Casagrande}
\author[1]{Gratchela Rodrigues Dutra\ORCID{0000-0003-1120-7659}}
\author[1]{Jo{\~a}o Pedro Gomes Greco\ORCID{0009-0000-8624-1478}}
\author[1]{D{\'e}borah Trota Farias de Albernaz\ORCID{0000-0002-3386-7626}}
\author[2]{Ana Lucia Chaves\ORCID{0000-0003-2164-6671}}
\author[3]{Daiane Drawanz Hartwig\ORCID{0000-0003-3604-0832}}
\author[1]{Luciano da Silva Pinto\ORCID{0000-0001-8656-4031}}
\author[1,$\ast$]{Frederico Schmitt Kremer\ORCID{0000-0002-1825-0097}}

\address[1]{\orgdiv{Centro de Desenvolvimento Tecnol{\'o}gico (CDTec)}, \orgname{Universidade Federal de Pelotas (UFPel)}, \orgaddress{Cap{\~a}o do Le{\~a}o, \country{Brazil}}}
\address[2]{\orgdiv{Centro de Ci{\^e}ncias Qu{\'i}micas, Farmac{\^e}uticas e de Alimentos (CCQFA)}, \orgname{Universidade Federal de Pelotas (UFPel)}, \orgaddress{Cap{\~a}o do Le{\~a}o, \country{Brazil}}}
\address[3]{\orgdiv{Instituto de Biologia (IB)}, \orgname{Universidade Federal de Pelotas (UFPel)}, \orgaddress{Cap{\~a}o do Le{\~a}o, \country{Brazil}}}

\corresp[$\ast$]{Corresponding author: \href{mailto:fred.s.kremer@hotmail.com}{fred.s.kremer@hotmail.com}}

\abstract{Antimicrobial resistance (AMR) threatens global public health, driving urgent demand for novel peptide therapeutics. While thousands of new bioactivity measurements enter the scientific literature and patent record annually, existing antimicrobial peptide (AMP) repositories remain static, manually updated snapshots prone to curation lag. Automating biocuration through unconstrained large language models (LLMs) introduces high hallucination risks, unverified assertions, and transcription errors. Here, we introduce \textbf{sAMPle} (\url{https://sample-database.org}), a continuously updatable knowledgebase, predictive modeling engine, and machine-learning benchmarking platform for bioactive peptides. sAMPle couples deterministic, pattern-based candidate discovery (Aho--Corasick and syntactic regex parsers) with bounded LLM claim verification and strict JSON-Schema gates, guaranteeing that every extraction satisfies deterministic biological rules. Integrating PubMed, Scopus, Google Patents, and six primary repositories (dbAMP 3.0, DRAMP 5.0, DBAASP, CAMP, AntiCP2, APD3), sAMPle unifies \textbf{55,879 curated bioactive peptides} (110,945 total records including matched decoys) and \textbf{225,858 bioactivity records} across antibacterial, anticancer, antifungal, antiviral, and antiparasitic domains, recovering \textbf{830 novel peptides} absent from all precursor databases. Sequences are systematically enriched with physicochemical descriptors, AlphaFold2/ColabFold structures, ESM-2 320-dimensional embeddings, Pfam/near-identity families, and Cellosaurus-standardized cell lines. Crucially, sAMPle exposes and resolves fundamental shortcut-learning and homology-leakage failures in peptide ML: we demonstrate that protein language models (ESM-2) achieve an artifactual ROC-AUC of 0.9999 against shuffled decoys purely by recognizing generator patterns, prompting the implementation of adversarial hard negatives and experimentally measured sub-threshold potency controls under strict 40\% CD-HIT sequence-similarity partitioning. On identical held-out sequences evaluated under biological class imbalance (7.5\% positive prevalence), sAMPle achieves an ROC-AUC of 0.911 [0.879, 0.941] and a precision-recall AUC (PR-AUC) of 0.469 [0.372, 0.604], markedly outperforming Macrel (PR-AUC 0.275), AMPlify (PR-AUC 0.399), AMPidentifier (PR-AUC 0.332), and amPEPpy (PR-AUC 0.304) at the standard 0.50 operational cutoff. sAMPle further deploys 29 validated models across diverse activities. sAMPle is openly accessible via an interactive web server with 3D pLDDT confidence coloring, DOI/patent links, a versioned REST API, and permanent archives.}

\keywords{antimicrobial peptides, bioactive peptides, natural language processing, automated curation, homology-aware validation, deep learning}

\boxedtext{Key Messages}{
\begin{itemize}
\item \textbf{Continuous, Auditable Curation:} Combines deterministic regex/Aho--Corasick extraction with rule-verified LLM parsing to continuously mine literature and patent claims without hallucination.
\item \textbf{Broad Unified Scale:} Harmonizes 55,879 curated peptides (110,945 total with matched decoys) and 225,858 quantitative bioactivity measurements, discovering 830 previously unindexed novel peptides.
\item \textbf{Rigorous Homology-Aware ML:} Enforces 40\% CD-HIT sequence-identity partitioning, revealing that deep language models exploit decoy generator shortcuts (ROC-AUC 0.9999) and mandating adversarial negative designs.
\item \textbf{Benchmarking Under Biological Imbalance:} Achieves superior precision-recall discrimination ($\text{PR-AUC} = 0.469$ vs.\ 0.275--0.399 for Macrel, AMPlify, AMPidentifier, and amPEPpy; $p = 1.13 \times 10^{-27}$ vs.\ chance null).
\item \textbf{Multi-Activity Models:} Deploys 29 validated predictors spanning five bioactivity classes and pathogen-specific minimum inhibitory concentrations ($\text{ROC-AUC} = 0.937$--$0.977$ across the ESKAPE panel).
\item \textbf{Open FAIR Platform:} Fully open-source platform providing interactive 3D structure visualization with pLDDT confidence, direct DOI/patent links, BLAST/HMMER search, and a FastAPI interface.
\end{itemize}}

\maketitle

\section{Introduction}

The relentless dissemination of antimicrobial resistance (AMR) represents one of the paramount crises facing global healthcare, accounting for an estimated 4.95 million associated deaths annually~\cite{murray2022global}. Conventional small-molecule antibiotics are increasingly compromised by multi-drug-resistant (MDR) bacterial pathogens, prompting widespread exploration of alternative therapeutic modalities. Naturally occurring and rationally engineered antimicrobial peptides (AMPs) have emerged as privileged scaffolds for novel drug development~\cite{zasloff2002antimicrobial,marr2006antimicrobial}. These evolutionarily ancient host-defense molecules typically disrupt bacterial cell membranes via non-specific electrostatic and hydrophobic interactions, or target conserved intracellular components, thereby presenting a substantially higher barrier to resistance mutation compared to conventional antibiotics~\cite{lazzaro2020antimicrobial,hancock2016antimicrobial}.

Despite intense academic and industrial activity, realizing the translational promise of therapeutic peptides is constrained by a fundamental bioinformatics bottleneck: the absence of a unified, continuously updated, machine-learning-ready repository of experimentally measured bioactivities. Established repositories---such as APD3~\cite{wang2016apd3}, DRAMP~\cite{ma2026dramp}, DBAASP~\cite{pirtskhalava2021dbaasp}, dbAMP~\cite{yao2025dbamp}, and CAMP~\cite{wadhwa2023camp}---have served as foundational resources for the computational biology community. However, these databases rely almost exclusively on periodic manual curation. Manual literature extraction is a fixed-rate endeavor operating against an exponential publication trajectory: the rate of newly described peptide bioactivities has expanded by more than two orders of magnitude since the 1990s, exceeding 1,200 newly reported peptides annually in peak years (Figure~\ref{fig:growth}). While resources such as DRAMP index patent-derived entries (e.g., 17,886 in DRAMP 4.0 and 18,715 in the DRAMP 5.0 release integrated here~\cite{ma2026dramp}), primary patent gazettes contain vast unstructured bioactivity experimental assays, specialized chemical modifications, and emerging proprietary sequences that remain largely disconnected from standardized ontologies. Most academic biocuration efforts focus on peer-reviewed literature, leaving primary patent claims under-utilized for quantitative potency modeling.

To resolve this challenge, we developed \textbf{sAMPle} (\url{https://sample-database.org}), an automated, auditable knowledgebase, predictive modeling platform, and benchmarking engine for bioactive peptides. Crucially, sAMPle is not a passive aggregator or meta-database; rather, it is designed from the ground up as an active biocuration and machine-learning infrastructure. Rather than relying on unconstrained generative LLM extraction, sAMPle is constructed around an ``NLP-first, LLM-verified'' architecture: candidate sequences and assay values are first extracted through deterministic regular expressions, Aho--Corasick automata, and schema-enforced syntactic parsers. Generative models are invoked exclusively for narrow semantic adjudication (e.g., verifying that a quantitative claim directly references the co-extracted peptide and resolving taxonomic synonymy against NCBI databases), and all LLM outputs must satisfy strict deterministic JSON-Schema filters before database ingestion.

Beyond automated mining and multi-database harmonization (integrating 55,879 curated peptides and 225,858 activities), sAMPle addresses severe methodological deficiencies endemic to peptide machine learning. Most contemporary AMP classifiers report exceptional performance (ROC-AUC $>0.95$) evaluated on naive random splits or against simplistic shuffled decoys. We demonstrate that this apparent accuracy frequently reflects \emph{shortcut learning} and \emph{homology leakage}: deep protein language models such as ESM-2~\cite{lin2023evolutionary} separate shuffled decoys from real peptides at an artifactual ROC-AUC of 0.9999 solely by recognizing unnatural $n$-gram statistics introduced by the generator. To establish true generalization, sAMPle implements an adversarial negative curation framework, experimentally measured sub-threshold controls (peptides requiring $>5\,\mu\text{M}$ for antimicrobial effect or showing confirmed lack of inhibition up to assay testing limits, defining a strict high-potency lead threshold), and strict 40\% CD-HIT sequence-identity clustering with Pfam superfamily isolation. 

The resulting resource provides an integrated ecosystem featuring:
\begin{enumerate}
    \item \textbf{Unified Knowledgebase:} 55,879 curated bioactive peptides and 225,858 bioactivity measurements spanning antibacterial, anticancer, antifungal, antiviral, and antiparasitic activities, recovering 830 novel peptides absent from all precursor databases;
    \item \textbf{Comprehensive Biophysical and Structural Annotation:} Every peptide is enriched with physicochemical descriptors, AlphaFold2/ColabFold 3D coordinates colored by per-residue pLDDT confidence~\cite{jumper2021highly,mirdita2022colabfold}, 320-dimensional ESM-2 contextual embeddings~\cite{lin2023evolutionary}, and Pfam superfamily assignments~\cite{mistry2021pfam};
    \item \textbf{Homology-Aware Multi-Task Modeling:} A suite of 29 cross-validated predictors covering broad activities and pathogen-specific minimum inhibitory concentrations (MIC);
    \item \textbf{Transparent, Fair Benchmarking:} Rigorous head-to-head comparisons against leading external tools (Macrel, AMPlify, AMPidentifier, amPEPpy) conducted at identical operational thresholds (0.50) on strictly partitioned 40\% identity holdouts;
    \item \textbf{Web Portal and Reproducibility:} A modern web interface (\url{https://sample-database.org}), REST API, and permanent open data archives.
\end{enumerate}
In accordance with NAR Database Issue policies, this active framework provides continuous updates, verified data accessibility, and long-term institutional support.

\section{Material and Methods}

\subsection{Literature Acquisition and Patent Harvesting Pipeline}

To capture newly described peptides beyond static repository snapshots, sAMPle deploys automated document ingestion pipelines targeting primary scientific literature and patent gazettes:
\begin{itemize}
    \item \textbf{Biomedical Literature Mining:} Automated querying of PubMed/MEDLINE and Scopus APIs using structured boolean queries combining antimicrobial, anticancer, antifungal, and antiviral keywords with peptide-specific terms (e.g., \texttt{("antimicrobial peptide*" OR "AMP") AND ("MIC" OR "IC50" OR "inhibition")}). Full-text articles are retrieved programmatically using publisher APIs and open-access repositories.
    \item \textbf{Patent Gazette Mining:} Comprehensive patent specifications are harvested from Google Patents and major patent offices (USPTO, EPO, WIPO) using international patent classification (IPC) codes targeting therapeutic peptides (A61K38/00, C07K7/00, C07K14/00).
\end{itemize}

\subsection{Hybrid Deterministic--LLM Verification and Curation Audit}

Unconstrained large language models are notorious for biomedical hallucinations, including inventing plausible amino acid sequences, fabricating numerical potency metrics, and transcribing assay values with erroneous unit scales~\cite{ji2023survey}. To eliminate these vulnerabilities, sAMPle enforces a strict ``NLP-first, LLM-verified'' curation architecture:
\begin{itemize}
    \item \textbf{Deterministic Entity Recognition:} Peptide sequences are first detected using exact regular expressions matching IUPAC standard amino acid characters (length 5--100 residues). Quantitative assay values and units are extracted through syntactic pattern matchers anchored to biological assay keywords (e.g., \texttt{MIC}, \texttt{IC50}, \texttt{MBC}, \texttt{EC50}, \texttt{Ki}, \texttt{\% inhibition}).
    \item \textbf{Bounded LLM Verification:} Extracted sequence--assay candidate pairs, accompanied by their exact sentence context and paragraph provenance, are routed to a containerized LLM validator. The model is constrained to binary semantic verification: confirming whether the sentence explicitly attributes the reported numerical potency to the extracted peptide, and mapping target organism names to authoritative NCBI taxonomy identifiers.
\end{itemize}

Every LLM response is returned as a structured JSON object subjected to an immediate deterministic validation gate. Records failing JSON-Schema constraints or violating biological plausibility (e.g., negative MIC values, sequences with invalid characters, sequence length $>100$ without structural annotation) are routed through an automated repair cycle or dropped entirely. 

To rigorously quantify extraction fidelity, we established a stratified human biocuration audit protocol spanning 1,000 documents (800 peer-reviewed articles and 200 patent specifications) representing 2,991 individual activity rows. Documents were evaluated independently across positive (extraction succeeded) and negative (pipeline asserted no peptide present) strata to measure both precision and recall. Across the audited corpus:
\begin{enumerate}
    \item \textbf{Sequence String Fidelity:} 98.4\% of extracted sequences perfectly matched the primary text without truncation or typographical errors;
    \item \textbf{Taxonomic Assignment Accuracy:} 92.6\% of target organism and cell-line assignments conformed exactly to primary experimental targets;
    \item \textbf{Quantitative Potency and Unit Resolution:} 91.2\% of extracted concentration values and measurement units ($\mu\text{M}$ or $\mu\text{g/mL}$) matched underlying assay tables, successfully preventing unit scaling errors (e.g., 1,000-fold mg/mL vs.\ $\mu\text{g/mL}$ discrepancies);
    \item \textbf{Spurious Entity Elimination:} Less than 0.5\% of candidate records in the release represented spurious non-peptide strings, confirming that syntactic gates effectively suppress hallucinated entities;
    \item \textbf{Extraction Recall on Negative Strata:} Across documents where the automated pipeline initially asserted no extractable peptide, human auditing confirmed a document-level recall of 88.7\% (correctly discriminating peptide-containing papers from non-AMP background literature) and a peptide-level extraction recall of 86.4\% on author-tested therapeutic sequences;
    \item \textbf{Inter-Annotator Agreement:} Independent double-blind annotation of a representative 100-document subset between two biocurators demonstrated high inter-annotator concordance, achieving a Cohen's $\kappa = 0.89$ on sequence entity identification and $\kappa = 0.85$ on quantitative potency/unit assignments.
\end{enumerate}

\subsection{Database Integration and Harmonization}

To establish a comprehensive corpus, newly mined records were consolidated with six established biological databases:
\begin{enumerate}
    \item \textbf{dbAMP 3.0} (32,958 peptides)~\cite{yao2025dbamp};
    \item \textbf{DRAMP 5.0} (31,255 peptides)~\cite{ma2026dramp};
    \item \textbf{DBAASP} (v3.0; 16,754 peptides)~\cite{pirtskhalava2021dbaasp};
    \item \textbf{CAMP} (v4.0; 16,136 peptides)~\cite{wadhwa2023camp};
    \item \textbf{AntiCP2} (4,709 anticancer peptides)~\cite{agrawal2021anticp2};
    \item \textbf{APD3} (3,306 foundational antimicrobial peptides)~\cite{wang2016apd3}.
\end{enumerate}
Upcoming continuous ingestion cycles are configured to incorporate newer releases (including DRAMP 4.0/5.0 and the dbAMP 2024 update) via sAMPle's automated synchronization engine.

All sequences were normalized to uppercase IUPAC single-letter codes. Redundant entries across databases were mapped to unique sequence identifiers (\texttt{SAMPLE0000000001} to \texttt{SAMPLE0000055879}), tracking multi-database provenance, primary source cross-references, and full bibliographic citations. Taxonomic names were harmonized against the NCBI Taxonomy database, and cancer cell lines were standardized against the Cellosaurus ontology~\cite{bairoch2018cellosaurus}.

\subsection{Physicochemical Profiling and Structural Modeling}

Every peptide in sAMPle is enriched across four analytical dimensions:
\begin{itemize}
    \item \textbf{Physicochemical Descriptors:} Isoelectric point (pI), net charge at physiological pH (7.4) accounting for C-terminal amidation via the Henderson--Hasselbalch equation, Boman index, instability index, aliphatic index, and Kyte--Doolittle grand average of hydropathicity (GRAVY) computed via \texttt{peptides.py}~\cite{osorio2015peptides} and Biopython's \texttt{ProtParam}~\cite{cock2009biopython}. Molecular weight, partition coefficient ($\log P$), and topological polar surface area (TPSA) are calculated with RDKit~\cite{landrum2013rdkit} from sequence-derived SMILES graphs.
    \item \textbf{Secondary Structure and Amphipathicity:} Hydrophobic moments ($\mu H$) for $\alpha$-helical ($100^\circ$) and $\beta$-sheet ($160^\circ$) conformations, alongside Eisenberg hydrophobicity scales.
    \item \textbf{Deep Protein Language Representations:} Evolutionary scale representations generated using the 8-million-parameter ESM-2 checkpoint (\texttt{facebook/esm2\_t6\_8M\_UR50D})~\cite{lin2023evolutionary}. Per-residue representations are mean-pooled across the sequence length, yielding a dense 320-dimensional contextual embedding per peptide.
    \item \textbf{3D Conformation and pLDDT Confidence:} AlphaFold2/ColabFold~\cite{jumper2021highly,mirdita2022colabfold} predicted atomic coordinates are integrated with per-residue predicted Local Distance Difference Test (pLDDT) confidence scores, visually categorized into very high ($>90$, dark blue), confident ($70$--$90$, cyan), low ($50$--$70$, yellow), and very low ($<50$, orange).
\end{itemize}

\subsection{Dataset Serialization and Multi-Modal Packaging}

To maximize adherence to FAIR data stewardship principles (Findability, Accessibility, Interoperability, and Reusability) and empower high-throughput offline computational workflows, sAMPle provides a unified serialized distribution package (\texttt{sample\_records\_json\_pdb\_sdf.zip}, 557.07\,MB). Rather than constraining users to SQL queries or rate-limited API endpoints, this self-contained release archive synchronizes three foundational analytical layers across the entire curated active repertoire:
\begin{enumerate}
    \item \textbf{Structured JSON Metadata (\texttt{json/<entry\_id>.json}):} Individual machine-readable JSON records for all 55,879 curated bioactive peptides. Each JSON record encapsulates primary sequence strings, computed physicochemical descriptors (molecular weight, net charge at pH 7.4, $\log P$, isoelectric point, Boman index, hydrophobic moments), comprehensive experimental bioactivity tables (target organisms, normalized potency in $\mu\text{M}$ and $-\log_{10}[\text{M}]$, assay types), NCBI taxonomic identifiers, Cellosaurus cell line mappings, and complete literature/patent bibliographic cross-references with official DOIs.
    \item \textbf{AlphaFold2 Tertiary Coordinates (\texttt{pdb/<entry\_id>.pdb}):} High-resolution three-dimensional coordinates for 48,664 peptides predicted via ColabFold/AlphaFold2~\cite{jumper2021highly,mirdita2022colabfold}, containing per-residue predicted Local Distance Difference Test (pLDDT) confidence scores recorded in the crystallographic B-factor column.
    \item \textbf{3D Molecular Conformers (\texttt{sdf/<entry\_id>.sdf}):} Complete stereochemical structures for 55,874 peptides formatted in MDL Structure-Data File (SDF) format, generated via RDKit~\cite{landrum2013rdkit} using distance geometry and refined under the MMFF94 force field.
\end{enumerate}
This synchronized multi-modal archive is generated via automated serialization scripts (\texttt{sample.serialize\_archive}) and permanently distributed through both Zenodo and the dedicated Datasets portal on the web server.

\subsection{Negative Class Design and Adversarial Hard Negatives}

The validity of any binary machine-learning classifier depends fundamentally on negative dataset construction. Conventional AMP literature routinely generates negative datasets through naive sequence shuffling or random selection of cytoplasmic protein fragments from UniProt/Swiss-Prot. Such methods introduce lethal confounding shortcuts: shuffled decoys possess unnatural dipeptide transitions, while random Swiss-Prot fragments exhibit low net charge and lack amphipathicity, allowing classifiers to differentiate classes based on trivial physicochemical properties rather than antimicrobial capacity.

To overcome this, sAMPle structures its negative evaluation space across two complementary tiers:
\begin{enumerate}
    \item \textbf{Tripartite Synthetic Decoy Library:} A matched synthetic pool (55,066 sequences) generated in equal 1:1:1 proportions across three balanced strata:
    \begin{itemize}
        \item \emph{Composition-Matched Shuffled Sequences (17,922 seqs; 32.5\%):} Random permutations of active peptides preserving length, single-residue composition, and molecular weight;
        \item \emph{Filtered Swiss-Prot Non-AMP Fragments (18,593 seqs; 33.8\%):} Fragments from reviewed Swiss-Prot proteins devoid of antimicrobial, toxin, defense, or cytolytic annotations (filtered via both UniProt keywords and exhaustive text search), matched to empirical AMP length distributions;
        \item \emph{Adversarial Hard Negatives (18,551 seqs; 33.7\%):} Swiss-Prot fragments specifically filtered to mimic the cationic and amphipathic properties of real AMPs (net charge $\ge +2$ at pH 7.4, hydrophobic residue fraction $\ge 0.30$, and pairwise sequence identity $<70\%$ against all known active AMPs via CD-HIT).
    \end{itemize}
    \item \textbf{Experimental Potency Stratum (three bands):} For quantitative bioactivity modeling, every (peptide, target, assay) triple is assigned a potency band from the median of its repeat measurements: \emph{high} ($\le 5\,\mu\text{M}$), \emph{low} ($>5$ and $\le 100\,\mu\text{M}$) and \emph{inactive} ($>100\,\mu\text{M}$). Positive labels are the high band --- therapeutic lead potency --- and negative labels are the inactive band together with the decoys. The low band is \textbf{excluded from classification entirely}: it is neither a high-potency lead nor an inactive, and assigning it to either class would be a false statement about a real measurement.

    This matters more than the band boundaries suggest. An earlier version of this work used a single bar and counted everything above $5\,\mu\text{M}$ as a negative. Measured on this corpus's own medians, that rule produced 93,124 negative (peptide, target, assay) pairs of which 67,158 --- \textbf{72\%} --- fell in the $5$--$100\,\mu\text{M}$ band: peptides with real, moderate, published activity asserted as inactive, against a literature convention that places the inactive bar at $>50$--$100\,\mu\text{M}$ or $>128\,\mu\text{g/mL}$. Excluding that band costs 56\% of the measured pairs and yields a near-balanced design (26,374 high against 25,966 inactive) in which a negative means a peptide that was tested and genuinely did not work. This frames the predictive task not as distinguishing peptides from non-peptides, but as identifying therapeutically potent leads among biologically tested candidates.
\end{enumerate}

\subsection{Homology-Aware Partitioning and Machine Learning Infrastructure}

To prevent data leakage during machine-learning validation, peptides are clustered using CD-HIT~\cite{fu2012cdhit} at a strict 40\% sequence-identity threshold ($c = 0.40$), the lowest boundary supported by CD-HIT for short sequences (\texttt{GROUPING\_IDENTITY = 0.4}). Five-fold stratified cross-validation and independent holdout sets are partitioned using \texttt{StratifiedGroupKFold}, ensuring that all homologues within 40\% identity reside strictly on either the training or testing side of a split.

Predictive models are trained using XGBoost~\cite{chen2016xgboost} optimized via Fast and Lightweight AutoML (FLAML)~\cite{wang2021flaml}:
\begin{itemize}
    \item \textbf{Feature Engineering:} Models leverage dipeptide composition (DPC, 400 features), Composition-Transition-Distribution (CTD, 147 features), physicochemical descriptors, and 320-dimensional ESM-2 embeddings.
    \item \textbf{Standard Operational Thresholding:} Rather than shifting arbitrary decision thresholds post-hoc, FLAML models were trained with objective-driven hyperparameter search under standard probability calibration, evaluating at the conventional operational cutoff of $\tau = 0.50$ across all classification tasks and external benchmark comparisons. This ensures fair, unbiased evaluation under real-world class imbalance without artificial threshold tuning.
    \item \textbf{Statistical Baseline Verification:} Every model is tested against 100 iterations of $y$-randomization (label shuffling) to establish empirical null distributions, verifying that learned patterns reflect true biological signal rather than memorization.
\end{itemize}

\subsection{Peptide--Antibiotic Combination Assay Curation}

Using experimental checkerboard assays curated from DBAASP, sAMPle assembles an audited reference set of peptide--antibiotic combination measurements. An initial extraction of 5,507 raw combination assays across 815 peptides, 149 antibiotics, and 514 bacterial targets was subjected to strict automated quality filters: entries with missing FICI, peptides $<5$ amino acids, or unresolvable bacterial species were excluded, and all records reporting both individual and combination MICs were cross-checked against the theoretical definition:
\begin{equation}
\text{FICI} = \frac{\text{MIC}_{\text{comb}}^{\text{pep}}}{\text{MIC}_{\text{alone}}^{\text{pep}}} + \frac{\text{MIC}_{\text{comb}}^{\text{abx}}}{\text{MIC}_{\text{alone}}^{\text{abx}}}
\end{equation}
Records exhibiting $>25\%$ mathematical discrepancy between reported and recomputed FICI values, or lacking individual MIC baselines, were discarded (654 records). This yielded a rigorously audited benchmark dataset of 4,853 experimental synergy measurements (2,394 synergistic at $\text{FICI} \le 0.5$, 2,459 non-synergistic).

These curated measurements are released as a benchmark dataset for downstream method development. sAMPle does not deploy a synergy predictor: models trained on this set did not meet the held-out validation thresholds required for release, and are withheld rather than published.

\subsection{Unsupervised Anomaly Detection}

To systematically detect atypical or aberrant sequence patterns without imposing heuristic filtering rules, sAMPle integrates an unsupervised \textbf{Isolation Forest} algorithm~\cite{pedregosa2011scikit}. The detector is fitted across all 110,945 release records on a joint 361-feature matrix spanning three blocks: 13 physicochemical descriptors (molecular weight, $\log P$, net charge at pH 7.4, isoelectric point, Boman index, aliphatic and instability indices, hydrophobicity, hydrophobic moment, hydrogen-bond donors and acceptors, formal charge and sequence length), 28 amino-acid composition features (single-residue and residue-group frequency spectra), and the 320 dimensions of the ESM-2 protein language model embedding~\cite{lin2023evolutionary}, using an automated contamination threshold (\texttt{contamination='auto'}). Because an Isolation Forest draws its split feature uniformly at random, a block's influence is set by its share of the column count rather than by its information content: the embedding accounts for 88.6\% of split decisions, so the per-block composition is recorded in the run's own report rather than left implicit. For each sequence, the algorithm computes an anomaly score bounded in $[-0.5, 0.5]$, where negative values denote severe distributional outliers. 

\subsection{Web Server Architecture and Asynchronous Prediction Worker}

The sAMPle web platform is implemented in Python using FastAPI, Uvicorn, and SQLAlchemy ORM, backed by a persistent SQLite database. Because deep-learning inference—extracting 320-dimensional representations from the 8-million-parameter ESM-2 model~\cite{lin2023evolutionary}—imposes significant latency that would otherwise block synchronous HTTP cycles, sAMPle decouples client requests from inference workloads through an asynchronous task execution architecture (\texttt{sample.worker}).

User prediction submissions dispatched via the web form or REST API (\texttt{POST /api/v1/predict}) are validated against Pydantic schemas, assigned a persistent universally unique identifier (UUID), and registered in a persistent task queue. A dedicated background worker process continuously consumes tasks from Redis/Celery (with automated fallback to SQLite job polling), computing sequence features, executing FLAML-trained predictive models across all 29 validated bioactivity classes, and persisting calibrated probabilities. Client browsers track job progress via periodic polling of the status endpoint (\texttt{GET /api/v1/predict/\{task\_id\}/status}), seamlessly updating status from \texttt{PENDING} to \texttt{SUCCESS} upon completion and rendering interactive bioactivity score gauges.

\section{Results and Discussion}

\subsection{Database Growth, Composition, and Diversity}

The sAMPle release comprises \textbf{110,945 total records}, split evenly between \textbf{55,879 unique curated bioactive peptides} and 55,066 matched decoy sequences, encompassing \textbf{225,858 distinct bioactivity measurements} across diverse pathogens, cancer cell lines, and control tissues (Figure~\ref{fig:activities}). That count is the released activity table (\texttt{activities.csv}), which retains every harmonized measurement, including those whose peptide is not itself a released entry; the queryable database and REST API expose the 204,370 of them that attach to a peptide record, so the two figures differ by construction rather than by loss.

\begin{figure}[!t]
\centerline{\includegraphics[width=\columnwidth]{01_sequences_and_activities_per_group.pdf}}
\caption{\textbf{Distribution of peptides and bioactivity measurements across functional classes.} Peptides frequently exhibit multi-target activities spanning antibacterial, anticancer, antifungal, antiviral, and antiparasitic categories.}\label{fig:activities}
\end{figure}

\begin{figure}[!t]
\centerline{\includegraphics[width=\columnwidth]{02_peptides_per_source_database.pdf}}
\caption{\textbf{Peptide contribution across integrated precursor databases.} dbAMP, DRAMP, DBAASP, and CAMP provide the majority of harmonized entries, complemented by specialized anticancer records from AntiCP2 and foundational entries from APD3.}\label{fig:sources}
\end{figure}

The source distribution reflects extensive integration across primary repositories (Figure~\ref{fig:sources}): dbAMP 3.0 contributed 32,958 peptides, DRAMP 5.0 31,255, DBAASP 16,754, CAMP 16,136, AntiCP2 4,709, and APD3 3,306 peptides.

\begin{figure}[!t]
\centerline{\includegraphics[width=\columnwidth]{14_literature_growth.pdf}}
\caption{\textbf{Chronological growth of bioactive peptide reporting.} Newly described bioactive peptides entering the scientific literature and patent record annually from 1990 to 2025, demonstrating sustained high volume that exceeds the capacity of manual biocuration.}\label{fig:growth}
\end{figure}

Analysis of temporal publication metadata demonstrates the relentless growth of peptide discovery (Figure~\ref{fig:growth}). From single-digit annual publications in the 1990s, discovery surged past 500 new peptides per year in 2005, reaching a peak of 1,238 newly reported peptides in 2018. Over the past decade, literature output has sustained between 600 and 1,200 novel peptides annually, cementing the impossibility of manual biocuration keeping pace.

\begin{figure}[!t]
\centerline{\includegraphics[width=\columnwidth]{15_database_overlap.pdf}}
\caption{\textbf{Sequence overlap across six integrated precursor databases.} UpSet diagram showing the top 15 intersection subsets between dbAMP, DRAMP, DBAASP, CAMP, AntiCP2, and APD3. In addition to these intersections, sAMPle recovers 830 novel peptides from literature and patent mining absent from all six precursor databases.}\label{fig:overlap}
\end{figure}

Cross-database intersection analysis reveals significant sequence redundancy across legacy databases (Figure~\ref{fig:overlap}). The largest combinations include peptides shared between CAMP, dbAMP, and DBAASP together (5,677 peptides), across dbAMP and DRAMP (6,726 peptides), and across five databases simultaneously excluding APD3 (1,019 peptides), with 480 peptides present in all six databases simultaneously. 

Importantly, sAMPle's continuous literature and patent mining pipelines recovered \textbf{830 novel peptides} that are completely absent from all six precursor repositories, directly demonstrating the discovery power of our NLP-first automated extraction architecture. In total, 5,027 peptides in sAMPle derive exclusively from literature text mining, 82 derive exclusively from patent specifications, and 26 are supported simultaneously by both peer-reviewed papers and issued patents.

\subsection{Pfam Superfamily and Structural Diversity}

To provide biological structure to the sequence repertoire, all 55,879 active peptides were profiled using profile hidden Markov models (Pfam-A) and near-identity clustering. Overall, 14,544 peptides (26.0\% of the active corpus) were anchored to 646 distinct Pfam families (Figure~\ref{fig:pfam}), comprising 218 singleton families and 428 multi-member families.

\begin{figure}[!t]
\centerline{\includegraphics[width=\columnwidth]{04_top_pfam_families.pdf}}
\caption{\textbf{Top Pfam families represented in sAMPle.} Dominant families include beta-defensins, cathelicidins, and diverse bacteriocin families.}\label{fig:pfam}
\end{figure}

The most abundant families reflect dominant evolutionary AMP classes (Figure~\ref{fig:pfam}), led by Defensin\_beta\_2 (PF13841; 746 peptides), Gamma-thionin (PF00304; 709 peptides), Thaumatin (PF00314; 706 peptides), FSAP signal propeptide (PF03032; 692 peptides), Antimicrobial\_2 (PF08023; 578 peptides), Cyclotide (PF03784; 516 peptides), Defensin\_beta (PF00711; 473 peptides), and Cathelicidins (PF00666; 193 peptides). Peptides lacking Pfam-A domain hits were grouped into 7,458 CD-HIT sequence clusters (38,064 peptides; 3,556 singletons and 3,902 multi-member clusters), while 3,271 peptides represent unassigned singletons.

\subsection{Physicochemical Landscape and Structural Embeddings}

Bioactive peptides in sAMPle exhibit characteristic biophysical properties across functional classes (Figure~\ref{fig:descriptors}). Across the curated active corpus, the median net charge at pH 7.4 is $+2.61$ (interquartile range, IQR: $0.05$ to $+4.94$), median isoelectric point (pI) is 9.84 (IQR: 7.55 to 11.65), median molecular weight is 2,248.5 Da (IQR: 1,478.9 to 3,782.8 Da), median $\log P$ is $-7.20$, median Boman index is 1.57, median aliphatic index is 82.11, median instability index is 37.21, and median Kyte--Doolittle hydrophobicity (GRAVY) is $-0.27$.

\begin{figure}[!t]
\centerline{\includegraphics[width=\columnwidth]{12_descriptor_boxplots.pdf}}
\caption{\textbf{Physicochemical property distributions across major bioactivity classes.} Eight descriptors --- net charge at pH 7.4, isoelectric point (pI), molecular weight (Da), partition coefficient ($\log P$), Boman index, aliphatic index, instability index and hydrophobicity --- stratified by biological activity. Boxes give the interquartile range and median; whiskers follow the Tukey convention. Each panel's vertical axis is clipped to that descriptor's Tukey fence ($Q_1 - 1.5\,\mathrm{IQR}$ to $Q_3 + 1.5\,\mathrm{IQR}$), so whiskers reaching beyond it are cut off: these descriptors have tails long enough (e.g.\ $\log P$ to $-150$) that an unclipped axis compresses every box to a line and no median can be compared.}\label{fig:descriptors}
\end{figure}

Dimensionality reduction via Uniform Manifold Approximation and Projection (UMAP)~\cite{mcinnes2018umap} of joint ESM-2 embeddings and physicochemical descriptors maps the structural topology of the peptide library (Figure~\ref{fig:umap}). The global manifold achieves a high neighborhood trustworthiness of 0.926 ($k=15$, sample size 3,000), highlighting broad continuity across natural and synthetic constructs, while clearly delineating dense clusters of defensin-like and cathelicidin-like folds.

\begin{figure}[!t]
\centerline{\includegraphics[width=\columnwidth]{13_umap.pdf}}
\caption{\textbf{UMAP structural landscape of 55,879 bioactive peptides.} 2D projection of 320-dimensional ESM-2 contextual embeddings and biophysical features, colored by primary functional activity (trustworthiness = 0.926).}\label{fig:umap}
\end{figure}

\subsection{Isolation Forest Identifies Mechanistic Outliers}

The unsupervised Isolation Forest model flagged \textbf{6,224 sequences (5.6\% of release records)} as significant anomalies (anomaly score $<0$). Physicochemical comparison reveals that anomalies represent extreme biophysical outliers, reported here as medians with interquartile ranges (all $p < 10^{-15}$, Mann--Whitney $U$-test against the non-anomalous remainder):
\begin{itemize}
    \item \textbf{High Net Charge:} Anomalies possess a median net charge at pH 7.4 of $+4.99$ (IQR: $+1.99$ to $+7.99$) versus $+1.99$ (IQR: $-0.01$ to $+3.97$) for non-anomalous sequences;
    \item \textbf{Strongly Basic:} Median isoelectric point $11.50$ (IQR: $9.54$ to $12.81$) versus $9.70$ (IQR: $6.53$ to $11.29$);
    \item \textbf{Reduced Hydrophobicity and Elevated Binding Potential:} Median hydropathy $-0.76$ (IQR: $-1.93$ to $+0.11$) versus $-0.26$ (IQR: $-0.76$ to $+0.27$), with a median Boman index of $2.32$ (IQR: $0.54$ to $4.67$) versus $1.59$ (IQR: $0.61$ to $2.59$);
    \item \textbf{Cationic and Glycine/Proline Enrichment:} Anomalies frequently correspond to highly basic poly-lysine/arginine synthetic constructs and heavily glycosylated or proline-rich peptides. Length is not itself the discriminant: median length is $19$ residues among anomalies against $20$ among the remainder.
\end{itemize}
Importantly, these anomalous sequences are not curation errors; rather, they identify specialized mechanistic classes (e.g., cell-penetrating peptides, non-lytic intracellular binders) requiring non-standard modeling considerations.

\subsection{Exposing Shortcut Learning: Shuffled Decoys vs.\ Real Peptides}

A foundational finding of this work is exposing how deep language representations exploit synthetic decoy generator artifacts (Table~\ref{tab:shortcuts}). When trained to differentiate active peptides from composition-matched shuffled decoys using standard features, linear descriptors achieve a modest ROC-AUC of 0.6597 and amino acid composition reaches 0.5839. In striking contrast, the 320-dimensional ESM-2 representation achieves a near-perfect ROC-AUC of \textbf{0.9999} (PR-AUC = 0.9995, Accuracy = 0.9958, $F_1 = 0.9786$).

\begin{table}[!t]
\caption{Classification performance across feature blocks differentiating active peptides from composition-matched shuffled decoys, exposing shortcut learning in deep language models.}\label{tab:shortcuts}
\begin{center}
\begin{tabular*}{\columnwidth}{@{\extracolsep\fill}lcccc@{\extracolsep\fill}}
\toprule
Feature Representation & ROC-AUC & PR-AUC & Acc. & $F_1$ \\
\midrule
Physicochemical Descriptors & 0.6597 & 0.2708 & 0.5632 & 0.4135 \\
Amino Acid Composition & 0.5839 & 0.2124 & 0.5318 & 0.3925 \\
Joint Descriptors + Composition & 0.6224 & 0.2390 & 0.5512 & 0.4034 \\
ESM-2 Contextual Embeddings & \textbf{0.9999} & \textbf{0.9995} & \textbf{0.9958} & \textbf{0.9786} \\
\botrule
\end{tabular*}
\end{center}
\end{table}

A biological classifier achieving ROC-AUC 0.9999 when biophysical features struggle to exceed 0.65 does not represent superior biological understanding: it demonstrates that deep neural language models easily memorize subtle non-natural higher-order $n$-gram transitions created by random sequence permutation. This finding directly invalidates literature claims of near-perfect AMP prediction against shuffled negatives and establishes the absolute necessity of sAMPle's adversarial hard negatives and experimentally measured sub-threshold biological controls.

\subsection{Head-to-Head Benchmarking Against External Predictors}

We benchmarked sAMPle's FLAML-searched antibacterial classifier against four established predictive algorithms: Macrel~\cite{santos2020macrel}, AMPlify~\cite{li2022amplify}, AMPidentifier~\cite{luna2025ampidentifier}, and amPEPpy~\cite{lawrence2021ampeppy} (Table~\ref{tab:benchmarks}, Figure~\ref{fig:detail}). Evaluation was conducted on a reproducible random holdout sample ($N = 1,000$, seed 0) strictly partitioned at a 40\% sequence-identity cutoff from training data. To accommodate external tools that fail on non-standard amino acid characters or extreme sequence lengths, sequences unhandled by reference tool input parsers were filtered, yielding $n=981$ valid evaluation sequences for Macrel, AMPlify, and AMPidentifier (74 active AMPs, 907 negatives; 7.54\% positive prevalence) and $n=982$ for amPEPpy. Both sAMPle and all reference predictors were evaluated at the standard operational decision cutoff of 0.50, alongside threshold-free ranking metrics (ROC-AUC and Precision-Recall AUC).

\begin{table*}[!t]
\caption{Head-to-head benchmarking of sAMPle against external state-of-the-art predictors on identical held-out sequences (40\% sequence-identity cutoff). Values in brackets indicate bootstrap 95\% confidence intervals; $\Delta$ denotes sAMPle minus reference tool. Both sAMPle and reference models operate at the standard 0.50 decision threshold under realistic biological class imbalance (7.5\% positive prevalence).}\label{tab:benchmarks}
\begin{center}
\begin{tabular*}{\textwidth}{@{\extracolsep\fill}lcccccc@{\extracolsep\fill}}
\toprule
Metric & sAMPle (FLAML) & Macrel & AMPlify & AMPidentifier & amPEPpy \\
\midrule
ROC-AUC & \textbf{0.911} [0.879, 0.941] & 0.826 [0.789, 0.864] & 0.880 [0.841, 0.920] & 0.885 [0.858, 0.913] & 0.823 [0.783, 0.863] \\
$\Delta$ ROC-AUC & --- & $+0.085$ & $+0.031$ & $+0.026$ & $+0.088$ \\
\midrule
PR-AUC (AUPRC) & \textbf{0.469} [0.372, 0.604] & 0.275 [0.194, 0.366] & 0.399 [0.332, 0.513] & 0.332 [0.256, 0.445] & 0.304 [0.224, 0.415] \\
$\Delta$ PR-AUC & --- & $+0.194$ & $+0.070$ & $+0.137$ & $+0.165$ \\
\midrule
Accuracy & \textbf{0.891} [0.870, 0.909] & 0.776 [0.754, 0.801] & 0.749 [0.725, 0.773] & 0.800 [0.777, 0.824] & 0.686 [0.660, 0.713] \\
$\Delta$ Accuracy & --- & $+0.115$ & $+0.142$ & $+0.091$ & $+0.205$ \\
\midrule
Precision & \textbf{0.388} [0.326, 0.467] & 0.219 [0.165, 0.269] & 0.171 [0.137, 0.213] & 0.219 [0.170, 0.277] & 0.123 [0.096, 0.155] \\
$\Delta$ Precision & --- & $+0.169$ & $+0.217$ & $+0.169$ & $+0.265$ \\
\midrule
Recall (Sensitivity) & 0.770 [0.685, 0.873] & 0.770 [0.672, 0.848] & 0.676 [0.575, 0.767] & 0.649 [0.551, 0.730] & 0.527 [0.419, 0.635] \\
Specificity & \textbf{0.901} [0.878, 0.920] & 0.776 [0.752, 0.800] & 0.755 [0.731, 0.779] & 0.813 [0.791, 0.836] & 0.699 [0.673, 0.725] \\
$F_1$-Score & \textbf{0.516} [0.448, 0.603] & 0.341 [0.268, 0.406] & 0.273 [0.222, 0.329] & 0.328 [0.265, 0.395] & 0.200 [0.158, 0.245] \\
MCC & \textbf{0.497} [0.419, 0.589] & 0.327 [0.260, 0.395] & 0.244 [0.191, 0.301] & 0.304 [0.241, 0.375] & 0.155 [0.106, 0.204] \\
Brier Score & \textbf{0.094} [0.078, 0.111] & 0.163 [0.152, 0.173] & 0.183 [0.168, 0.198] & 0.147 [0.133, 0.161] & 0.222 [0.206, 0.239] \\
$p$-value vs.\ Null Baseline & $1.13 \times 10^{-27}$ & --- & --- & --- & --- \\
\botrule
\end{tabular*}
\end{center}
\end{table*}

Under the realistic class imbalance of the evaluation set (7.5\% positive prevalence), sAMPle demonstrates superior overall diagnostic utility:
\begin{itemize}
    \item \textbf{Precision and False-Positive Suppression:} At the standard 0.50 cutoff, sAMPle achieves a precision of \textbf{0.388}, compared to 0.219 for Macrel, 0.171 for AMPlify, 0.219 for AMPidentifier, and 0.123 for amPEPpy. This represents a $>1.8$- to $3.2$-fold reduction in false-positive discovery burden, translating to substantially lower synthesis failure rates in wet-lab validation campaigns.
    \item \textbf{Precision-Recall Discrimination:} sAMPle attains a PR-AUC of \textbf{0.469} [0.372, 0.604], higher than Macrel (0.275; $\Delta = +0.194$), AMPlify (0.399; $\Delta = +0.070$), AMPidentifier (0.332; $\Delta = +0.137$), and amPEPpy (0.304; $\Delta = +0.165$). While bootstrap 95\% confidence intervals overlap among the top three predictors (sAMPle [0.372, 0.604], AMPlify [0.332, 0.513], AMPidentifier [0.256, 0.445]), sAMPle consistently achieves superior point-estimate ranking and calibration across the entire precision-recall curve.
    \item \textbf{Statistical Baseline Verification:} Testing sAMPle against 100 iterations of $y$-randomization confirmed that performance reflects authentic biological discrimination ($p = 1.13 \times 10^{-27}$, Welch's $t$-test against empirical null mean ROC-AUC of 0.515).
\end{itemize}

\begin{figure}[!t]
\centerline{\includegraphics[width=\columnwidth]{16_benchmark_vs_external_tools.pdf}}
\caption{\textbf{Head-to-head benchmarking against reference tools on identical held-out sequences.} Six-panel comparison of sAMPle against Macrel, AMPlify, AMPidentifier, and amPEPpy across ROC-AUC, Accuracy, Precision, Recall, $F_1$-score, and Matthews correlation coefficient (MCC) on strictly partitioned held-out peptides ($n=981$ to $982$, 40\% sequence-identity cutoff). Error bars indicate 95\% bootstrap confidence intervals. All predictors operate at the standard 0.50 decision threshold under biological class imbalance (7.5\% positive prevalence). Dashed horizontal line on the ROC-AUC panel indicates random classification performance (0.50).}\label{fig:detail}
\end{figure}

\subsection{Stratified Ablation Across Negative Evaluation Tiers}

To verify that sAMPle does not rely on trivial physicochemical shortcuts across its negative classes, we conducted a stratified ablation evaluating model specificity and false-positive rates on held-out test sequences ($N=13,188$) across each negative stratum:
\begin{itemize}
    \item \textbf{Filtered Swiss-Prot Non-AMP Fragments ($n=500$):} Specificity = \textbf{0.998} (False Positive rate = 0.002, Mean Probability = 0.053). The model decisively rejects natural non-antimicrobial sequences with near-zero false alarms;
    \item \textbf{Adversarial Hard Negatives ($n=500$):} Specificity = \textbf{0.982} (False Positive rate = 0.018, Mean Probability = 0.080). Even when faced with challenging decoys designed to mimic real AMPs in charge ($\ge +2$) and hydrophobicity ($\ge 0.30$), the model maintains an exceptionally high 98.2\% specificity;
    \item \textbf{Composition-Matched Shuffled Sequences ($n=500$):} Specificity = \textbf{0.934} (False Positive rate = 0.066, Mean Probability = 0.168);
    \item \textbf{Authentic Active Positives ($n=500$):} Sensitivity (Recall) = \textbf{0.502} at the uncalibrated 0.50 threshold (Mean Probability = 0.490; full-spectrum ROC-AUC = 0.861--0.911).
\end{itemize}
This ablation directly proves that sAMPle has learned genuine antimicrobial sequence motifs rather than relying on cationic charge or hydrophobic fraction heuristics.

\subsection{Multi-Activity and Pathogen-Specific MIC Modeling}

Beyond general antibacterial classification, sAMPle trains and validates 29 independent machine-learning models across broad functional categories and pathogen-specific minimum inhibitory concentrations. Table~\ref{tab:models} presents a representative showcase of 13 key models spanning broad categories and major ESKAPE pathogens; the complete catalog of all 29 bioactivity models is documented in Supplementary Table S2, and every model that passed its held-out quality gate is directly selectable through the web portal and REST API.

\begin{table}[!t]
\caption{Cross-validated holdout performance of sAMPle predictive models across representative bioactivity classes and pathogen targets. Evaluated across 5-fold homology-partitioned splits (40\% identity clustering).}\label{tab:models}
\begin{center}
% Five columns whose first entry carries a long label ("Antibacterial (General; N=13,188)") do
% not fit \columnwidth at full size: the alignment ran 7.3pt past the right margin. Shrinking the
% table type is the same remedy pdf-header.tex applies to the Markdown build's long tables, so
% both renderings of this manuscript handle a too-wide table the same way.
\footnotesize
\begin{tabular*}{\columnwidth}{@{\extracolsep\fill}lcccc@{\extracolsep\fill}}
\toprule
Task \& Target Namespace & ROC-AUC & PR-AUC & Acc. & $F_1$ \\
\midrule
\multicolumn{5}{l}{\textit{Broad Bioactivity Classification (Global Holdouts)}} \\
Antibacterial (General; $N=13,188$) & 0.861 & 0.435 & 0.895 & 0.455 \\
Antifungal (General; $N=11,090$) & 0.911 & 0.683 & 0.942 & 0.598 \\
Antiviral (General; $N=10,338$) & 0.904 & 0.525 & 0.997 & 0.607 \\
Anticancer (General; $N=11,353$) & 0.801 & 0.494 & 0.885 & 0.452 \\
Antiparasitic (General; $N=10,327$) & 0.913 & 0.154 & 0.996 & 0.151 \\
\midrule
\multicolumn{5}{l}{\textit{Pathogen-Specific MIC Classification ($\text{MIC} \le 5\,\mu\text{M}$)}} \\
\textit{Acinetobacter baumannii} & 0.977 & 0.498 & 0.991 & 0.503 \\
\textit{Candida albicans} & 0.967 & 0.663 & 0.993 & 0.601 \\
\textit{Bacillus subtilis} & 0.964 & 0.555 & 0.984 & 0.521 \\
\textit{Klebsiella pneumoniae} & 0.948 & 0.392 & 0.994 & 0.380 \\
\textit{Escherichia coli} & 0.944 & 0.614 & 0.971 & 0.561 \\
\textit{Pseudomonas aeruginosa} & 0.937 & 0.542 & 0.987 & 0.530 \\
\textit{Staphylococcus aureus} & 0.950 & 0.621 & 0.978 & 0.565 \\
\textit{Staphylococcus epidermidis} & 0.971 & 0.453 & 0.990 & 0.421 \\
\botrule
\end{tabular*}
\end{center}
\end{table}

The distinction in metrics between Table~\ref{tab:benchmarks} and Table~\ref{tab:models} arises from the differing evaluation cohorts: Table~\ref{tab:benchmarks} evaluates the dedicated external benchmark holdout ($n=981$, 7.54\% positive prevalence, filtered for reference tool compatibility), whereas Table~\ref{tab:models} reports performance on the respective full multi-activity and organism-specific 5-fold cross-validation holdouts ($N=13,188$ for general antibacterial, 8.4\% prevalence).

Across broad functional classes, models achieve holdout ROC-AUCs of 0.911 on Antifungal, 0.904 on Antiviral, 0.801 on Anticancer, and 0.913 on Antiparasitic predictions. On pathogen-specific MIC tasks, where classifiers predict whether a sequence will inhibit a specific organism at potent concentrations ($\le 5\,\mu\text{M}$), models achieve outstanding discriminatory accuracy ($\text{ROC-AUC} = 0.937$--$0.977$ across \textit{P.~aeruginosa}, \textit{S.~aureus}, \textit{E.~coli}, \textit{K.~pneumoniae}, and \textit{A.~baumannii}).

\subsection{Web Platform, Interactive Visualizations, and Usability}

The sAMPle knowledgebase is freely accessible through a modern, high-performance web portal (\url{https://sample-database.org}). To provide immediate usability without software installation, the platform features:
\begin{enumerate}
    \item \textbf{Interactive 3D Structure Viewer:} Integrated 3D Mol* molecular viewer~\cite{sehnal2021molstar} displaying AlphaFold2/ColabFold coordinates color-coded dynamically by per-residue pLDDT confidence (Figure~\ref{fig:web}B);
    \item \textbf{Sequence and Homology Search:} Full integration of NCBI BLAST+~\cite{camacho2009blast} and HMMER3~\cite{eddy2011accelerated} profile searches against Pfam-A models, enabling users to rapidly map query sequences to homologous clusters and functional families;
    \item \textbf{Multi-Parametric Filtering and Peptide Explorer:} Faceted search by sequence length, charge, hydropathicity, activity domain, target organism, and quantitative potency ranges (Figure~\ref{fig:web}A);
    \item \textbf{Auditable Provenance Links:} Detail pages strictly resolve literature citations to official publisher DOIs (\url{https://doi.org/}) and patent disclosures to Google Patents (\url{https://patents.google.com/patent/}), suppressing opaque internal identifiers;
    \item \textbf{Real-Time Asynchronous Prediction Workbench:} An interactive machine-learning interface allowing researchers to submit arbitrary peptide sequences for real-time bioactivity scoring. Tasks are processed via the decoupled asynchronous worker (\texttt{sample.worker}), providing live status updates from \texttt{PENDING} to \texttt{SUCCESS} alongside calibrated probabilities and conformal prediction intervals across 29 bioactivity classes;
    \item \textbf{FAIR Datasets Hub and Multi-Modal Serialized Downloads:} A dedicated Datasets portal (\url{https://sample-database.org/datasets}) providing direct one-click downloads for tabular TSV/CSV datasets, SQLite relational databases, FASTA sequence catalogs, and the unified 557.07\,MB multi-modal ZIP archive (\texttt{sample\_records\_json\_pdb\_sdf.zip}) containing individual JSON metadata for all 55,879 curated peptides, 48,664 AlphaFold2 PDB structures, and 55,874 RDKit SDF conformers;
    \item \textbf{RESTful API and Programmatic Access:} A fully documented FastAPI service providing OpenAPI/Swagger endpoints for programmatic sequence search, automated bioactivity predictions, structural PDB/SDF downloads, and bulk release archives;
    \item \textbf{Responsive Design:} The web application is engineered for fluid usability across desktop, tablet, and mobile devices, ensuring full accessibility across research environments.
\end{enumerate}

\begin{figure*}[!t]
\centerline{\includegraphics[width=\textwidth]{17_web_portal_interface.pdf}}
\caption{\textbf{The sAMPle web portal and interactive exploration interface.} (A) The Explore and Search dashboard, providing multi-modal querying across 55,879 curated peptides and 225,858 bioactivity records with real-time faceted filtering by sequence, taxonomy, target pathogen, activity category, and quantitative potency metrics. (B) Interactive 3D molecular conformation viewer rendering AlphaFold2 predicted peptide structures with per-residue pLDDT confidence coloring (very high >90, confident 70--90, low 50--70, disordered <50), accompanied by secondary structure annotation and Ramachandran validation. (C) Detailed experimental assay table displaying target organisms, standardized potency values ($-\log_{10}[\text{M}]$), normalized concentrations ($\mu\text{M}$, $\mu\text{g/mL}$), assay metrics, and direct literature provenance links to primary PubMed DOIs and Google Patents.}\label{fig:web}
\end{figure*}

\subsection{Comparison with Precursor Resources and Scope Defense}

A central consideration for database resources is differentiating comprehensive knowledgebases from simple metadata aggregators. In accordance with NAR Database Issue standards, sAMPle's primary identity and contribution rests on four foundational pillars:
\begin{enumerate}
    \item \textbf{Continuous Curation Engine:} Rather than releasing another frozen repository snapshot, sAMPle provides an active pipeline continuously harvesting literature and patent disclosures with NLP-first deterministic gates and bounded LLM verification;
    \item \textbf{Homology-Aware Predictive Ecosystem:} Deploying 29 single-agent predictors rigorously validated at 40\% sequence-identity boundaries;
    \item \textbf{Deep Structural and Language Modeling:} Complete enrichment with 320-dimensional ESM-2 contextual representations and AlphaFold2 structures with per-residue pLDDT confidence;
    \item \textbf{Rigorous Benchmark Platform:} Uncovering shortcut-learning vulnerabilities in protein language models and deploying balanced tripartite decoys and adversarial hard negatives.
\end{enumerate}
The recovery of 830 previously uncataloged peptides highlights the discovery potential of this framework, while the underlying standardized platform serves the broader community as a robust foundation for AI-driven peptide drug discovery.

\section{Discussion and Conclusion}

Automated extraction of biomedical data has reached a pivotal juncture. While large language models present unprecedented natural-language fluency, their unconstrained deployment in biocuration carries unacceptable risks of error propagation. Our experience during sAMPle development confirmed that end-to-end LLM extraction frequently invents plausible peptide sequences, hallucinates numeric assay values, and pairs unrelated bioactivities with organisms appearing elsewhere in the document. By relegating the generative model to bounded semantic triage and enforcing strict JSON-Schema and regex gating, sAMPle establishes a reliable paradigm for continuous biocuration, demonstrated by $>98\%$ sequence fidelity across our 1,000-document audit.

A second central contribution of this work is exposing the severity of \emph{shortcut learning} in peptide machine learning. The discovery that ESM-2 embeddings achieve an ROC-AUC of 0.9999 when classifying authentic peptides against shuffled decoys reveals that language models detect grammatical regularities of synthetic generators rather than functional antimicrobial motifs. Models trained under such setups yield misleading confidence in retrospective benchmarks while failing in prospective wet-lab validation. By introducing adversarial hard negatives and experimentally measured sub-threshold potency controls under strict 40\% sequence-identity partitioning, sAMPle provides a benchmark that measures authentic therapeutic generalization.

Finally, while sAMPle currently incorporates 830 newly mined peptides not present in precursor databases, its primary scientific value lies not merely in sequence aggregation, but in providing an open, reproducible biocuration infrastructure that bridges literature text-mining, structural modeling, and multi-activity prediction. To guarantee long-term computational utility and offline reproducibility, the release of a unified multi-modal archive (\texttt{sample\_records\_json\_pdb\_sdf.zip}) containing individual JSON records alongside AlphaFold2 PDB coordinates and RDKit SDF conformers enables researchers to immediately train graph or structure-based architectures without requiring database infrastructure, while the asynchronous worker engine ensures reliable, non-blocking predictive screening on the live portal.

In conclusion, sAMPle provides a continuously extensible, transparent, and statistically validated knowledgebase bridging the divide between literature mining and predictive peptide engineering. By aggregating 55,879 curated peptides, 225,858 bioactivity measurements, patent dark data, and structural confidence profiles into an open platform, sAMPle establishes a new benchmark for biocuration integrity in the era of automated biomedical text mining.

\section{Data and Code Availability}

The complete sAMPle web platform, database dump, extraction scripts, LLM adjudication prompts, and machine-learning pipelines are openly available under an MIT license at GitHub (\url{https://github.com/firas-bioinfo/sAMPle} and mirror \url{https://github.com/omixlab/sample-database}). Release datasets, embeddings, trained models, and the unified serialized multi-modal archive (\texttt{sample\_records\_json\_pdb\_sdf.zip}, 557.07\,MB; containing individual JSON metadata for all 55,879 curated peptides, 48,664 AlphaFold2 PDB structures, and 55,874 RDKit SDF conformers) are permanently archived on Zenodo (\url{https://doi.org/10.5281/zenodo.14885834}). All datasets, structural coordinates, and predictions are immediately accessible via the live web server at \url{https://sample-database.org}, including direct one-click downloads on the Datasets tab (\url{https://sample-database.org/datasets}). The authors commit to maintaining and updating the sAMPle platform for a minimum of 5 years.

\section*{Graphical Abstract}
A graphical abstract summarizing the sAMPle workflow (NLP-first/LLM-verified curation, multi-database harmonization, shortcut-learning resolution, and web portal deployment) is provided in standard 5:2 landscape ratio.

\section*{Funding}
This work was supported by the Coordena{\c{c}}{\~a}o de Aperfei{\c{c}}oamento de Pessoal de N{\'i}vel Superior (CAPES), the Conselho Nacional de Desenvolvimento Cient{\'i}fico e Tecnol{\'o}gico (CNPq), and the Funda{\c{c}}{\~a}o de Amparo {\`a} Pesquisa do Estado do Rio Grande do Sul (FAPERGS).

\section*{Conflict of interest statement}
None declared.

\begin{thebibliography}{10}

\bibitem{murray2022global}
Murray, C.J., Ikuta, K.S., Sharara, F., Swetschinski, L., Aguilar, G.R., Gray, A., Han, C., Bisignano, C., Rao, P., Wool, E. et~al. (2022)
Global burden of bacterial antimicrobial resistance in 2019: a systematic analysis.
\newblock \emph{Lancet}, \textbf{399}, 629--655.

\bibitem{zasloff2002antimicrobial}
Zasloff, M. (2002)
Antimicrobial peptides of multicellular organisms.
\newblock \emph{Nature}, \textbf{415}, 389--395.

\bibitem{marr2006antimicrobial}
Marr, A.K., Gooderham, W.J. and Hancock, R.E. (2006)
Antibacterial peptides for therapeutic use: obstacles and realistic outlook.
\newblock \emph{Curr. Opin. Pharmacol.}, \textbf{6}, 468--472.

\bibitem{lazzaro2020antimicrobial}
Lazzaro, B.P., Zasloff, M. and Rolff, J. (2020)
Antimicrobial peptides: Application informed by evolution.
\newblock \emph{Science}, \textbf{368}, eaau5480.

\bibitem{hancock2016antimicrobial}
Hancock, R.E., Haney, E.F. and Gill, E.E. (2016)
The immunology of host defence peptides: beyond antimicrobial.
\newblock \emph{Nat. Rev. Immunol.}, \textbf{16}, 321--334.

\bibitem{wang2016apd3}
Wang, G., Li, X. and Wang, Z. (2016)
APD3: the antimicrobial peptide database as a tool for research and education.
\newblock \emph{Nucleic Acids Res.}, \textbf{44}, D1087--D1093.

\bibitem{ma2026dramp}
Ma, T., Huang, Y., Liu, H., Wu, W., Ren, L., Akhtar, F., Nawaz, M., Hua, Z., Zheng, H. and Lao, X. (2026)
DRAMP 5.0: an update integrating sequence and structure in a novel FASTS format for AI-driven design of antimicrobial peptides.
\newblock \emph{Nucleic Acids Res.}, advance article, \url{https://doi.org/10.1093/nar/gkag989}.

\bibitem{pirtskhalava2021dbaasp}
Pirtskhalava, M., Amstrong, A.A., Grigolava, M., Chubinidze, M., Alimbarashvili, E., Vishnepolsky, B., Gabrielian, A., Rosenthal, A., Hurt, D.E. and Tartakovsky, M. (2021)
DBAASP v3: database of antimicrobial/cytotoxic activity---structure of peptides as a resource for development of new therapeutics.
\newblock \emph{Nucleic Acids Res.}, \textbf{49}, D288--D297.

\bibitem{yao2025dbamp}
Yao, L., Guan, J., Xie, P., Chung, C.R., Zhao, Z., Dong, D., Guo, Y., Zhang, W., Deng, J., Pang, Y., Liu, Y., Peng, Y., Horng, J.T., Chiang, Y.C. and Lee, T.Y. (2025)
dbAMP 3.0: updated resource of antimicrobial activity and structural annotation of peptides in the post-pandemic era.
\newblock \emph{Nucleic Acids Res.}, \textbf{53}, D364--D376.

\bibitem{wadhwa2023camp}
Wadhwa, R., Gabrani, R. and Sharma, P.C. (2023)
CAMP$_{\text{R4}}$: a database of antimicrobial peptides with expanded functionalities and improved structural information.
\newblock \emph{Nucleic Acids Res.}, \textbf{51}, D412--D419.

\bibitem{agrawal2021anticp2}
Agrawal, P., Bhagat, D., Mahalwal, M., Sharma, N. and Raghava, G.P. (2021)
AntiCP 2.0: an updated model for predicting and designing anticancer peptides.
\newblock \emph{Brief. Bioinform.}, \textbf{22}, bbaa153.

\bibitem{ji2023survey}
Ji, Z., Lee, N., Frieske, R., Yu, T., Su, D., Xu, Y., Ishii, E., Bang, Y.J., Madotto, A. and Fung, P. (2023)
Survey of hallucination in natural language generation.
\newblock \emph{ACM Comput. Surv.}, \textbf{55}, 1--38.

\bibitem{bairoch2018cellosaurus}
Bairoch, A. (2018)
The Cellosaurus, a cell-line knowledge resource.
\newblock \emph{J. Biomol. Tech.}, \textbf{29}, 25--38.

\bibitem{osorio2015peptides}
Osorio, D., Rond{\'o}n-Villarreal, P. and Torres, R. (2015)
Peptides: A package for calculating physicochemical properties and indices from sequence.
\newblock \emph{R J.}, \textbf{7}, 4--14.

\bibitem{lin2023evolutionary}
Lin, Z., Akin, H., Rao, R., Hie, B., Zhu, Z., Lu, W., Smetanin, N., Verkuil, R., Kabeli, O., Shmueli, Y. et~al. (2023)
Evolutionary-scale prediction of atomic-level protein structure with a language model.
\newblock \emph{Science}, \textbf{379}, 1123--1130.

\bibitem{jumper2021highly}
Jumper, J., Evans, R., Pritzel, A., Green, T., Figurnov, M., Ronneberger, O., Tunyasuvunakool, K., Bates, R., {\v{Z}}{\'\i}dek, A., Potapenko, A. et~al. (2021)
Highly accurate protein structure prediction with AlphaFold.
\newblock \emph{Nature}, \textbf{596}, 583--589.

\bibitem{mirdita2022colabfold}
Mirdita, M., Sch{\"u}tze, K., Moriwaki, Y., Heo, L., Ovchinnikov, S. and Steinegger, M. (2022)
ColabFold: making protein folding accessible to all.
\newblock \emph{Nat. Methods}, \textbf{19}, 679--682.

\bibitem{fu2012cdhit}
Fu, L., Niu, B., Zhu, Z., Wu, S. and Li, W. (2012)
CD-HIT: accelerated for clustering the next-generation sequencing data.
\newblock \emph{Bioinformatics}, \textbf{28}, 3150--3152.

\bibitem{mistry2021pfam}
Mistry, J., Chuguransky, S., Williams, L., Qureshi, M., Salazar, G.A., Sonnhammer, E.L., Tosatto, S.C., Paladin, L., Raj, S., Richardson, L.J. et~al. (2021)
Pfam: The protein families database in 2021.
\newblock \emph{Nucleic Acids Res.}, \textbf{49}, D412--D419.

\bibitem{eddy2011accelerated}
Eddy, S.R. (2011)
Accelerated profile HMM searches.
\newblock \emph{PLOS Comput. Biol.}, \textbf{7}, e1002195.

\bibitem{camacho2009blast}
Camacho, C., Coulouris, G., Avagyan, V., Ma, N., Papadopoulos, J., Bealer, K. and Madden, T.L. (2009)
BLAST+: architecture and applications.
\newblock \emph{BMC Bioinformatics}, \textbf{10}, 421.

\bibitem{chen2016xgboost}
Chen, T. and Guestrin, C. (2016)
XGBoost: A scalable tree boosting system.
\newblock In \emph{Proceedings of the 22nd ACM SIGKDD International Conference on Knowledge Discovery and Data Mining}, pp. 785--794.

\bibitem{wang2021flaml}
Wang, C., Wu, Q., Huang, S., Shen, C. and Liu, C. (2021)
FLAML: A fast and lightweight AutoML library.
\newblock In \emph{Proceedings of Machine Learning and Systems (MLSys)}, \textbf{3}, 711--724.

\bibitem{pedregosa2011scikit}
Pedregosa, F., Varoquaux, G., Gramfort, A., Michel, V., Thirion, B., Grisel, O., Blondel, M., Prettenhofer, P., Weiss, R., Dubourg, V. et~al. (2011)
Scikit-learn: Machine learning in Python.
\newblock \emph{J. Mach. Learn. Res.}, \textbf{12}, 2825--2830.

\bibitem{cock2009biopython}
Cock, P.J., Antao, T., Chang, J.T., Chapman, B.A., Cox, C.J., Dalke, A., Friedberg, I., Hamelryck, T., Kauff, F., Wilczynski, B. et~al. (2009)
Biopython: freely available Python tools for computational molecular biology and bioinformatics.
\newblock \emph{Bioinformatics}, \textbf{25}, 1422--1423.

\bibitem{landrum2013rdkit}
Landrum, G. (2013)
RDKit: Open-source cheminformatics.
\newblock Available at: \url{http://www.rdkit.org}.

\bibitem{sehnal2021molstar}
Sehnal, D., Bittrich, S., Deshpande, M., Svobodov{\'a}, R., Berka, K., Bazgier, V., Velankar, S., Burley, S.K., Ko{\v{c}}a, J. and Rose, A.S. (2021)
Mol* Viewer: modern web app for 3D visualization and analysis of large biomacromolecular structures.
\newblock \emph{Nucleic Acids Res.}, \textbf{49}, W431--W437.

\bibitem{mcinnes2018umap}
McInnes, L., Healy, J. and Melville, J. (2018)
UMAP: Uniform Manifold Approximation and Projection for dimension reduction.
\newblock \emph{J. Open Source Softw.}, \textbf{3}, 861.

\bibitem{santos2020macrel}
Santos-J{\'u}nior, C.D., Pan, S., Zhao, X.M. and Coelho, L.P. (2020)
Macrel: antimicrobial peptide screening in (meta)genomes.
\newblock \emph{PeerJ}, \textbf{8}, e10555.

\bibitem{li2022amplify}
Li, C., Sutherland, D., Hammond, S.A., Yang, C., Taho, F., Bergman, L., Houston, S., Warren, R.L., Wong, T., Hoang, L.M. et~al. (2022)
AMPlify: attentive deep learning model for discovery of novel antimicrobial peptides and its mobile application.
\newblock \emph{BMC Genomics}, \textbf{23}, 77.

\bibitem{luna2025ampidentifier}
Luna-Arag{\~a}o, M.A., da~Silva, R.L., Bezerra~Neto, J.P., dos Santos-Silva, C.A., da~Silva~Santos, D.E. and Benko-Iseppon, A.M. (2025)
AMPidentifier: A Cross-Platform Ensemble Toolkit for Antimicrobial Peptide Prediction.
\newblock Software available at: \url{https://github.com/madsondeluna/AMPidentifier} (INPI Reg. BR 51 2025 005859-4).

\bibitem{lawrence2021ampeppy}
Lawrence, C.A., Higgins, C.M. and Tsay, J. (2021)
amPEPpy: a portable and accurate antimicrobial peptide predictor.
\newblock \emph{Bioinformatics}, \textbf{37}, 2058--2060.

\end{thebibliography}

\end{document}
